\textbf{Объект разработки:} процесс программной реализации алгоритмов эллиптической криптографии для экосистемы .NET.

\textbf{Цель работы:} проектирование и реализация высокопроизводительных, кроссплатформенных криптографических примитивов Curve25519 и Ed25519 на управляемом языке C\# с применением техник аппаратного ускорения и оптимизации работы с памятью.

\textbf{Методы и методология:} математическое моделирование арифметики конечных полей с использованием избыточного представления Radix-51; применение полных законов сложения (кривые Эдвардса) для обеспечения защиты от атак по времени; векторизация вычислений с применением SIMD-инструкций (AVX2); использование паттернов Zero Allocation (\texttt{Span<T>}, \texttt{ref struct}) для исключения накладных расходов на сборку мусора; бенчмаркинг.

\textbf{Результаты работы и их новизна:} разработана полностью нативная для платформы .NET библиотека, реализующая обмен ключами X25519 и электронную подпись Ed25519 (RFC 7748, RFC 8032) за строго константное время. Впервые для управляемой среды C\# реализована пакетная (batched) обработка криптографических операций с использованием инструкций AVX2. Практические результаты тестов показали превосходство разработанного решения над индустриальным стандартом Bouncy Castle в среднем на 30\% при массовой обработке данных.

\textbf{Область применения:} высоконагруженные серверные .NET-приложения, микросервисы, реализации протоколов TLS 1.3, IoT-шлюзы, где требуется обработка тысяч криптографических операций в секунду при сохранении кроссплатформенности без использования P/Invoke.